% Nouveau EGA V (cote 25): archive cote 25, site page 146 (PDF page 147).
% typescript. Draft compared: none.
% AI-transcribed by Claude Opus 5.5 (claude-opus-5-5, adaptive thinking, effort high), Message Batches API, batch ega25b (msgbatch_01A7n5YuXubUPNK1zsepi8RG), 2026-11-13; not yet checked by eye.
\scanpage{25/146}
\mspage{7}

p et q tels que $p+q=m+1$ , on utilise la construction précédente relati
vement aux morphismes $P\longrightarrow M/Q = S$ et $Q\longrightarrow M/P = R$ , pour trouver
des sections \xout{dx} a de $P \otimes \det S \otimes \det P^{\unsure{\hand{-}}1} = P \otimes \det M \otimes \det P^{-1} \otimes$
$\det Q^{-1}$ et b de $.. = Q \otimes \det M \otimes \det P^{-1} \otimes \det Q^{-1}$ , qu'on peut
aussi envisager comme des homomorphismes de $\underline{L} = \det P \otimes \det Q \otimes \det M^{-1}$
dans P resp. Q . (NB on dénote, pour un Module loc libre F, par dét F
sa puissance extérieure maxima, et on utilise que, pour une suite
exacte $0 \quad F' \quad F \quad F'' \quad 0$ de tels Modules, on a un isomorphisme
canonique $\det F \quad \det F' \otimes \det F''$ ). Ceci posé, \uline{on a commutativité}
\uline{dans le diagramme}

\diagram{
\begin{tikzcd}
 & P \arrow[r, hook] & M \\
\det P\otimes\det Q\otimes\det M^{-1}=L \arrow[r, phantom] & L \arrow[u, "a"] \arrow[r, "b"] & Q \arrow[u, hook]
\end{tikzcd}
}
,

(évdentuellement au signe près, cela dépend peut-être des conventions
adoptées pour définir certains des isomorphismes canoniques utilisés,
qui comportent le choix d'un signe ...), donc les idéaux $\unsure{\underline{I}}(a)$ et $\unsure{\underline{I}}(b)$
sont identiques (puisque P et Q sont loc facteurs directs dans M). \hand{NB en}
\hand{les points $\notin V(a)=V(b)$ , L est exactement l'intersection de P et Q dans M.}

Il reste, pour \unsure{clariffer|clarifier} complètement la situation particulière
envisagée avec X/Y/S et $\phi$ , décrite dans le diagramme \inserted{(D)} plus haut, \xout{dx}
d'expliciter a et b . On trouve d'abord
\[ \det P = \hand{\omega}_{X/Y}^{\unsure{\hand{-}}1} \otimes \underline{O}_{\unsure{X|V}}1 \quad , \quad \det Q = \hand{\underline{\omega}}_{Y/S} \otimes \underline{O}_{V}1 \]
\[ \det M = \hand{\omega}_{X/S} \otimes \underline{O}_{V}1 \ = \hand{\omega}_{X/Y} \otimes \hand{\omega}_{Y/S} \otimes \underline{O}_{V}1 \]
(le dernier isomorphisme provenant de la suite exacte de transitivité
pour les $\hand{\underline{\Omega}}^1$ pour les morphismes\xout{x} lisses $X-Y$ et $Y-S$), d'où
\[ \underline{L} \hand{\simeq} \hand{\omega}_{X/Y}{}^{-2} \otimes \underline{O}_{V}1 \]
(NB $\hand{\omega}_{X/Y}$ désigne \xout{lxxp} dét $\hand{\underline{\Omega}}^1_{X/Y}$ ) . Ceci dit, \uline{l'\unsure{h}omorphisme} \uline{a} ce \hand{---}
\uline{factorise par la composante $\underline{O}_{V1}$}

\mspage{146}
